\subsection{Main Single-Qubit Gates}

\subsubsection{Identity Gate (I)}

The \definitionWithSpecificIndex{Identity Gate}{Single Qubit Gates: Identity Gate}{}, denoted $I$, is the most elementary gate in quantum computing. It performs \textbf{no change} to the state of a qubit. Its action on a qubit is trivial, yet it \textbf{serves important roles in quantum circuits}, particularly \hl{when managing multi-qubit systems} or \hl{timing synchronization}.

\highspace
\begin{flushleft}
    \textcolor{Green3}{\faIcon{square-root-alt} \textbf{Matrix Representation}}
\end{flushleft}
The Identity gate is represented by the $2 \times 2$ \textbf{identity matrix}:
\begin{equation}
    I = \begin{pmatrix}
        1 & 0 \\ 0 & 1
    \end{pmatrix}
\end{equation}
\textbf{Applying} the Identity gate to any qubit \textbf{leaves the state unchanged}. Formally, for a qubit in state $\ket{\psi}$:
\begin{equation*}
    I \ket{\psi} = \ket{\psi}
\end{equation*}

\highspace
To better understand its effect, consider its action on the computational basis:
\begin{align*}
    I\ket{0} =& \begin{pmatrix}
        1 & 0 \\ 0 & 1
    \end{pmatrix}
    \begin{pmatrix}
        1 \\ 0
    \end{pmatrix} =
    \begin{pmatrix}
        1 \\ 0
    \end{pmatrix} = \ket{0} \\
    I\ket{1} =& \begin{pmatrix}
        1 & 0 \\ 0 & 1
    \end{pmatrix}
    \begin{pmatrix}
        0 \\ 1
    \end{pmatrix} =
    \begin{pmatrix}
        0 \\ 1
    \end{pmatrix} = \ket{1}
\end{align*}
Hence, the Identity gate \textbf{preserves the state} of the qubit, including any superposition or entangled state it may be in.

\highspace
\begin{flushleft}
    \textcolor{Green3}{\faIcon{book} \textbf{Physical Interpretation}}
\end{flushleft}
Although it may appear useless at first glance, the Identity gate has \textbf{conceptual and practical utility}:
\begin{itemize}[label=\textcolor{Green3}{\faIcon{check}}]
    \item It acts as a \textbf{placeholder} or \textbf{``\emph{do nothing}'' operation} when required by \textbf{circuit timing}.
    \item In \textbf{multi-qubit circuits}, it allows one \textbf{qubit to remain unchanged} while other qubits are operated on.
    \item In \textbf{matrix composition} of gates, it serves to \textbf{align dimensions} when performing \textbf{tensor products} (e.g., $I \otimes H$ applies Hadamard to the second qubit only).
\end{itemize}

\highspace
\begin{flushleft}
    \textcolor{Green3}{\faIcon{question-circle} \textbf{Bloch Sphere Perspective}}
\end{flushleft}
In the Bloch sphere (\autopageref{fig:bloch-sphere}), the \textbf{identity gate} represents a \textbf{rotation of zero degrees}, meaning the state vector remains exactly where it is. It is the neutral element of quantum transformations.
